% Exact candidate proof for source and production-Rethlas audit.
% This file is isolated from the live book. It makes no novelty or priority
% claim and must not be promoted before a final/correct production receipt.

\begin{proof}
Put \(\Gamma=-\log|K^\times|\).  We first record the three chart facts
used below.  A finite integral piecewise-linear tropicalization on a good
strict analytic space is locally covered by one smooth analytic moment map
and a piecewise integral-affine map; see
\cite[Proposition~8.3 and \S8.4]{GJR26}.  On an open subset of an algebraic
analytification, \cite[Proposition~7.2]{Gub16} replaces such a finite
analytic moment, after shrinking, by an algebraic moment with the same
tropicalization.  It is equality of the tropicalizations, rather than of
the analytic functions themselves, that will be used.  By factoring the
algebraic moment through the canonical moment map of a common very affine
open, as in \cite[\S4.12]{Gub16}, one obtains an ordinary algebraic
tropical chart dominating any finitely many prescribed charts.

The third fact is the full-ambient condition in the definition of an
ordinary delta-metric.  In an ordinary delta-chart, the facewise Hessian of
the weight is the restriction of a \((1,1)\)-superform on the full ambient
vector space with piecewise-smooth coefficients
\cite[Definition~9.9(ii) and \S9.13]{GK17}.  A piecewise-smooth coefficient
is one single-valued continuous function across the cells
\cite[Definition~1.6]{GK17}; equivalently, the ambient superform is a
codimension-zero delta-preform
\cite[Remark~1.7, Example~2.10, and Remark~9.10]{GK17}.  Thus its
coefficient tensor has one value at a wall, not one unrelated value for
each incident face.

We next establish the specialization step that will return from a product
Gauss point to the original variety.  Let \(Y\) be an algebraic
\(K\)-variety, let \(y\in Y^{\mathrm{an}}\), and write
\(H_y=\mathcal H(y)\).  For \(w\in N_{\mathbf R}\), let

\[
             \zeta_{w,y}\in T_{H_y}^{\mathrm{an}}
                    \subset (T\times Y)^{\mathrm{an}}
\]

be the relative Gauss point characterized by

\begin{equation}
\label{eq:p09p2-relative-gauss-point}
 \left|\sum_{m\in M}a_m\chi^m\right|_{\zeta_{w,y}}
   =\max_m |a_m|\exp(-\langle m,w\rangle),
 \qquad a_m\in H_y.
\end{equation}

This is a point of the analytic fibre product, not a naive ordered pair of
Berkovich points.  The Shilov-section construction of
\cite[\S\S12.1.10--12.1.11]{CLD25} also shows that

\[
             (w,y)\longmapsto\zeta_{w,y}
\]

is continuous.

We claim that, if \(w\in N\otimes_{\mathbf Z}\Gamma\), every ordinary
analytic neighbourhood \(W\) of \(\zeta_{w,y}\) contains \(i_t(y)\) for
some \(t\in T(K)\) satisfying \(\operatorname{trop}_T(t)=w\), where

\[
                       i_t:Y\longrightarrow T\times Y,
                       \qquad z\longmapsto(t,z)
\]

is the constant algebraic section.  Choose a basis of \(M\).  Since
\(w\in N\otimes\Gamma\), scaling the corresponding torus coordinates by
elements of \(K^\times\) reduces the assertion to the unit Gauss point of

\[
 \mathcal T_y=\mathcal M\bigl(
 H_y\langle Z_1^{\pm1},\ldots,Z_r^{\pm1}\rangle\bigr).
\]

The intersection \(W\cap\mathcal T_y\) contains a basic ordinary
neighbourhood of the Gauss point controlled by the absolute values of
finitely many analytic functions.  After normalizing each nonzero function
by a coefficient attaining its Gauss norm, its initial reduction is a
nonzero Laurent polynomial over \(\widetilde H_y\).  It is therefore enough
to choose one element of \((\widetilde K^\times)^r\) outside the zero loci
of finitely many nonzero Laurent polynomials with coefficients in
\(\widetilde H_y\).

For completeness, \((\widetilde K^\times)^r\) is Zariski dense in
\((\widetilde H_y^\times)^r\).  More generally, if \(k\subset L\) are
fields with \(k\) infinite, a nonzero Laurent polynomial over \(L\) cannot
vanish on all of \((k^\times)^r\).  This follows by induction on \(r\):
choose the first \(r-1\) entries so that one nonzero coefficient polynomial
in the last variable remains nonzero, and then avoid the finitely many
roots of the resulting one-variable polynomial.  Here the isometric
embedding \(K\hookrightarrow H_y\) induces
\(\widetilde K\hookrightarrow\widetilde H_y\), while \(\widetilde K\) is
algebraically closed and hence infinite.  Lifting the chosen residue tuple
to \((K^{\circ\times})^r\), the ultrametric inequality gives equality with
the Gauss norm for every one of the finitely many functions.  Multiplying
back the scaling coordinates produces the required \(t\in T(K)\).  This
proves the claim without invoking rigid-point density.

Consequently, let \(u\) be a real-valued function germ at \(y\) and let
\(m\in M_{\mathbf R}\).  If

\begin{equation}
\label{eq:p09p2-gauss-specialization}
 \operatorname{pr}_Y^*u
   +\langle m,\operatorname{trop}_T\circ\operatorname{pr}_T\rangle
\end{equation}

is CLD-smooth at \(\zeta_{w,y}\), choose a neighbourhood \(W\) on which
it is smooth and then choose \(t\) as above with \(i_t(y)\in W\).  The set
\(i_t^{-1}(W)\) is an ordinary neighbourhood of \(y\).  Pulling
\eqref{eq:p09p2-gauss-specialization} back by \(i_t\) shows there that

\[
 u+\langle m,\operatorname{trop}_T(t)\rangle
\]

is smooth.  The second term is constant, so \(u\) is smooth at \(y\).

Assume now that the metric with weight \(F\circ h\) is an ordinary
delta-metric.  Put

\[
 \rho=\operatorname{trop}_T\circ\operatorname{pr}_T
           +h\circ\operatorname{pr}_X
       :(T\times X)^{\mathrm{an}}\longrightarrow N_{\mathbf R}.
\]

The equivariance identity \eqref{eq:p09p2-product-equivariance} gives

\begin{equation}
\label{eq:p09p2-action-pullback-weight}
                         F\circ\rho
                    =(F\circ h)\circ\mu^{\mathrm{an}}.
\end{equation}

Algebraic pullback preserves ordinary delta-metrics by
\cite[Proposition~9.15]{GK17}.  Hence \(F\circ\rho\) is an ordinary
delta-weight on \(T\times X\).

Fix \(x\in X^{\mathrm{an}}\), choose provisionally
\(w_0\in N\otimes\Gamma\), and take an ordinary delta-chart at
\(\zeta_{w_0,x}\).  Cover all coordinates of the locally integral PL map
\(\rho\) by one analytic moment, algebraize that finite moment by
\cite[Proposition~7.2]{Gub16}, adjoin a basis of characters of the first
\(T\)-factor, and pass to one common ordinary algebraic refinement.  Since
the line bundle is trivial, we use its global frame \(1\); if the original
delta-presentation used another algebraic frame, adjoining their algebraic
unit ratio changes the weight only by an affine tropical coordinate and
does not change its facewise Hessian.  After shrinking, we obtain

\[
 \Phi:V\longrightarrow P\subset\mathbf R^s,
 \qquad
 H:P\longrightarrow N_{\mathbf R},
\]

where \(\Phi\) is the tropicalization in an ordinary algebraic chart,
\(H\) is piecewise integral affine, and

\begin{equation}
\label{eq:p09p2-common-chart}
 \rho|_V=H\circ\Phi,
 \qquad
 \mathrm d'_{\mathrm P}\mathrm d''_{\mathrm P}(F\circ H)
          =\gamma|_P
\end{equation}

for a \((1,1)\)-superform \(\gamma\) on \(\mathbf R^s\) with
piecewise-smooth coefficients.  The old ambient superform has been pulled
back along the integral-affine projection from the refined chart.

Fix common locally finite complexes of definition for \(H\), the tropical
support, and every coefficient function of \(\gamma\), and restrict to a
small target neighbourhood meeting only finitely many cells.  All chart
coordinates, factorizations, and walls have now been fixed; only now do we
choose the final Gauss weight.

For \(w'\) sufficiently close to \(w_0\), continuity of the Shilov section
keeps \(\zeta_{w',x}\) in \(V\).  Set

\[
                         z(w')=\Phi(\zeta_{w',x}).
\]

Each algebraic unit defining the final chart restricts on the fibre over
\(x\) to an element of \(\mathcal H(x)(M)^\times\), hence to a quotient of
two finite Laurent polynomials.  By \eqref{eq:p09p2-relative-gauss-point},
its logarithmic norm is the difference of the maxima of finitely many
integral-affine functions of \(w'\).  Thus \(z\) is finite piecewise
integral affine on a sufficiently small neighbourhood of \(w_0\).

Refine its finite source subdivision by the inverse images of all the
already fixed target cells.  Since \(K\) is algebraically closed,
\(\Gamma\) is nonzero and divisible; it is therefore dense in \(\mathbf R\).
We may choose

\[
                         w\in N\otimes_{\mathbf Z}\Gamma
\]

arbitrarily close to \(w_0\), in the relative interior of a
full-dimensional cell of the source refinement.  Put \(z_0=z(w)\).  There
are an open neighbourhood \(D\) of \(w\), an injective linear map

\begin{equation}
\label{eq:p09p2-gauss-slope-map}
                         a:N_{\mathbf R}\hookrightarrow\mathbf R^s,
\end{equation}

and a cell \(C\) of the common target refinement such that

\begin{equation}
\label{eq:p09p2-gauss-common-cell}
 z(w+u)=z_0+a(u),
 \qquad
 z_0\in\operatorname{relint}(C),
 \qquad
 z(D)\subset C.
\end{equation}

The character coordinates of the first torus factor make \(a\) injective.
Since \(D\) is full dimensional, \eqref{eq:p09p2-gauss-common-cell} gives

\begin{equation}
\label{eq:p09p2-gauss-lineality}
                         a(N_{\mathbf R})\subset L_C.
\end{equation}

Moreover, for \(u\) sufficiently small,

\[
 H(z(w+u))=\rho(\zeta_{w+u,x})=w+u+h(x).
\]

Differentiation on \(C\) gives \(D_CH\circ a=\operatorname{id}\).  If
\(\sigma\) is any cell in the star of \(C\), then \(C\) is a face of
\(\sigma\), the affine restrictions \(H|_C\) and \(H|_\sigma\) agree on
\(C\), and \(a(N_{\mathbf R})\subset L_C\subset L_\sigma\).  Hence

\begin{equation}
\label{eq:p09p2-H-lineality}
              D_\sigma H(a(v))=v
 \qquad
 (v\in N_{\mathbf R},\ \sigma\in\operatorname{Star}(C)).
\end{equation}

We now use the full-ambient tensor in
\eqref{eq:p09p2-common-chart}.  Let \(G_0\) be the coefficient bilinear
form of \(\gamma\) at \(z_0\).  On the relative interior of a cell
\(\sigma\) in the star of \(C\), the facewise Hessian of \(F\circ H\) is

\[
              b(D_\sigma H\,\mathord\cdot,
                D_\sigma H\,\mathord\cdot).
\]

The coefficient functions of \(\gamma\) have one common value at
\(z_0\).  Taking the limit to \(z_0\) through the relative interior of
\(\sigma\) therefore yields

\begin{equation}
\label{eq:p09p2-wall-tensor}
 \left.G_0\right|_{L_\sigma\times L_\sigma}
 =b(D_\sigma H\,\mathord\cdot,D_\sigma H\,\mathord\cdot).
\end{equation}

Fix \(m\in M\), choose \(v=(b^\sharp)^{-1}(m)\), where

\[
 b^\sharp:N_{\mathbf R}\longrightarrow M_{\mathbf R},
 \qquad v\longmapsto b(v,\mathord\cdot),
\]

and put

\[
                         \alpha_m=G_0(a(v),\mathord\cdot).
\]

By \eqref{eq:p09p2-gauss-lineality}, \eqref{eq:p09p2-H-lineality}, and
\eqref{eq:p09p2-wall-tensor},

\[
 \left.\alpha_m\right|_{L_\sigma}
   =b(v,D_\sigma H\,\mathord\cdot)
   =\left.\mathrm d\langle m,H\rangle\right|_\sigma.
\]

The function \(\langle m,H\rangle\) is continuous and affine on every
cell, and every cell in the star contains \(C\) as a face.  Its value at
\(z_0\) and its derivative on every incident cell therefore agree with one
ambient affine function.  On the whole star germ,

\begin{equation}
\label{eq:p09p2-affine-coordinate}
 \langle m,H(z)\rangle
  =\langle m,H(z_0)\rangle+\alpha_m(z-z_0).
\end{equation}

It follows that \(\langle m,\rho\rangle\) is CLD-smooth at
\(\zeta_{w,x}\).  Since

\[
 \langle m,\rho\rangle
 =\langle m,\operatorname{trop}_T\circ\operatorname{pr}_T\rangle
    +\operatorname{pr}_X^*\langle m,h\rangle,
\]

the specialization argument applied to
\(u=\langle m,h\rangle\) shows that this coordinate is smooth at \(x\).
The point \(x\) and the lattice coordinate \(m\) were arbitrary, so every
coordinate of \(h\) is CLD-smooth.

Conversely, suppose that the coordinates of \(h\) in a basis of \(M\) are
CLD-smooth.  Near any point, concatenate their finitely many analytic
moment presentations into one analytic moment.  The coordinate vector
\(h\), and hence \(F\circ h\), is the pullback of a smooth function on its
tropical target.  Apply \cite[Proposition~7.2]{Gub16} once to replace the
finite analytic moment by an algebraic moment with the same
tropicalization, and factor it through the canonical moment of a very
affine open.  Thus \(F\circ h\) is smooth in the ordinary algebraic-chart
sense.  Smooth metrics are ordinary delta-metrics by
\cite[Remark~9.16]{GK17}.  Real linear combinations show that testing one
basis of \(M\) is equivalent to testing all coordinates.  This proves the
first assertion.

Before applying it to Raynaud uniformization, we verify one locality fact
not contained in \cite[Proposition~9.15]{GK17}.  Let

\[
                 \iota:V_X\xrightarrow{\sim}V_Y
\]

be an analytic isomorphism between analytic open subsets of algebraic
\(K\)-varieties.  Suppose algebraic line bundles on the two varieties are
identified over these opens by an analytic isomorphism, and transport a
metric through this identification.  Then the metric is ordinary delta on
one side if and only if it is ordinary delta on the other.

Indeed, pull the finitely many algebraic moment coordinates of a
delta-presentation on \(V_Y\) back by \(\iota\), and adjoin the invertible
analytic function relating the pulled-back algebraic frame to an algebraic
frame on the source.  Package this finite family as one analytic moment.
By \cite[Proposition~7.2]{Gub16}, after shrinking it has the same
tropicalization as an algebraic moment.  Factoring through the canonical
moment map gives an ordinary algebraic chart.  On its target, the old
weight is composed with an integral-affine projection and is changed only
by the last affine coordinate coming from the frame ratio; the old ambient
codimension-zero superform pulls back along that projection.  This is an
ordinary delta-presentation on the source.  Applying the same argument to
\(\iota^{-1}\) proves the equivalence.  The assertion is only being made
for analytic opens of algebraic varieties and algebraic line bundles.

We now consider the split Raynaud data in the statement.  The Raynaud
extension \(E\) is algebraic and its \(T\)-action is algebraic by
\cite[\S16]{GJR26}; compare \cite[Remark~B.6.6]{CLD25}.  In the split
situation, the quotient

\[
                       p:E^{\mathrm{an}}\longrightarrow A^{\mathrm{an}}
\]

is an analytic covering by \(\Lambda\), hence locally an analytic
isomorphism; see \cite[\S19.2.1 and \S\S B.5.12--B.5.13]{CLD25}.  The
locality fact just proved applies to the algebraic line bundle on \(E\)
in the following structure formula.

By \cite[Proposition~19.2.5]{CLD25}, there are a line bundle \(M_L\) on
\(B\), a formal metric on \(q^*M_L\), an isomorphism of its analytification
with \(p^*L\), and a quadratic function \(F_L\) on \(N_{\mathbf R}\) such
that the pullback of the normalized cubical metric is the product of that
formal metric and the metric with weight

\begin{equation}
\label{eq:p09p2-raynaud-weight}
                              F_L\circ\tau.
\end{equation}

Let \(Q_L\) denote the homogeneous quadratic form of Raynaud's
polarization data, and let \(b_L\) be its associated symmetric bilinear
form, normalized by \(Q_L(w)=b_L(w,w)\).  The same proposition says that
the homogeneous degree-two part of \(F_L\) is \(\frac12Q_L\); hence

\begin{equation}
\label{eq:p09p2-raynaud-quadratic}
             F_L(w)=\frac12b_L(w,w)+\ell_L(w)+c_L.
\end{equation}

The canonical tropicalization

\[
                       \tau:E^{\mathrm{an}}\longrightarrow N_{\mathbf R}
\]

is locally integral PL and satisfies

\[
              \tau(t\mathbin{\cdot}e)
                 =\operatorname{trop}_T(t)+\tau(e).
\]

Since \(L\) is ample, \cite[Theorem~B.6.4]{CLD25} says that \(Q_L\) is
positive definite; thus \(b_L\) is nondegenerate.  Formal metrics are
piecewise linear by \cite[Proposition~8.11]{GK17}, hence ordinary
delta-metrics by \cite[Remark~9.16]{GK17}.  Ordinary delta-metricity is
local, invariant under isometry, and closed under tensor products and duals
by \cite[\S9.14]{GK17}.  The analytic-local equivalence above and
\eqref{eq:p09p2-raynaud-weight} therefore give

\[
 \overline L_{\mathrm{cub}}\text{ is ordinary delta on }A^{\mathrm{an}}
 \quad\Longleftrightarrow\quad
 F_L\circ\tau\text{ is ordinary delta on }E^{\mathrm{an}}.
\]

Apply the first assertion with \(X=E\), \(h=\tau\), and \(b=b_L\).
It follows that the two equivalent conditions in the preceding display
hold if and only if every coordinate of \(\tau\) is CLD-smooth.  Finally,
\cite[Remark~19.2.3]{CLD25} identifies simultaneous smoothness of these
coordinates with isotriviality of the residual extension of \(B_s\) by
\(T_{\widetilde K}\).  Every cubical metric on a fixed line bundle
is a positive constant multiple of the normalized one
\cite[Lemma~19.1.6(b)]{CLD25}; multiplication by a constant changes the
weight by a
smooth constant and preserves ordinary delta-metricity.  This proves the
claimed equivalence and its non-isotrivial mixed-reduction consequence.
\end{proof}
